\section{Стратегии управления}

\subsection{Стратегия 1: Стабилизация только угла}
Выберем управление в виде обратной связи по углу и угловой скорости:
\[
\ddot{u} = a\varphi + b\dot{\varphi}, \quad a,b \in \mathbb{R}. \tag{2}
\]

Подставляя (2) в (1), получаем замкнутую систему:
\[
\ddot{\varphi} + \frac{b}{l}\dot{\varphi} + \frac{a - g}{l}\varphi = 0. \tag{3}
\]

\subsubsection{Анализ устойчивости}
Характеристическое уравнение (3):
\[
\lambda^2 + \frac{b}{l}\lambda + \frac{a - g}{l} = 0.
\]

Для асимптотической устойчивости необходимо и достаточно:
\[
\frac{b}{l} > 0 \quad\Rightarrow\quad b > 0, \qquad \frac{a - g}{l} > 0 \quad\Rightarrow\quad a > g.
\]

\subsection{Стратегия 2: Одновременная стабилизация угла и точки опоры}
Для возврата точки опоры в нуль добавим обратную связь по \(u\) и \(\dot{u}\):
\[
\ddot{u} = -a\varphi - b\dot{\varphi} - cu - d\dot{u}, \quad a,b,c,d > 0. \tag{4}
\]

Система уравнений в переменных состояния:
\[
\begin{cases}
\dot{x}_1 = x_2, \\
\dot{x}_2 = \dfrac{g}{l}x_1 - \dfrac{1}{l}x_4, \\
\dot{x}_3 = x_4, \\
\dot{x}_4 = -a x_1 - b x_2 - c x_3 - d x_4,
\end{cases}
\]
где \(x_1 = \varphi\), \(x_2 = \dot{\varphi}\), \(x_3 = u\), \(x_4 = \dot{u}\).

\begin{figure}[H]
	\centering
	\includegraphics[width=0.7\linewidth]{03-block-diagram.png}
	\caption{Структурная схема системы управления}
	\label{fig:block}
\end{figure}